\section{Background}\label{background}
In this section, we review some necessary backgrounds which form the basis of our work, i.e., (1) Matrix Factorization (MF) in POI Recommendation, (2) decentralized learning.

\subsection{Matrix Factorization in POI Recommendation}

MF aims to learn user and item (POI) latent factors through regressing over the existing user-item ratings \cite{koren2009matrix,mnih2007probabilistic}, which can be formalized as follows,
\begin{equation}\small
\label{lfm}
\begin{split}
\mathop {\arg \min }\limits_{u_i,v_j} \sum\limits_{i=1}^{I} \sum\limits_{j=1}^{J} {{\left( {r_{ij} - u_i^T{v_j}} \right)}^2}  + \lambda\left(\sum\limits_{i=1}^{I} {||{u_i}||^2}  + \sum\limits_{j=1}^{J} {||{v_j}||^2}\right),
\end{split}	
\end{equation}
where $u_i$ and $v_j$ denote the latent factors of user $i$ and item $j$, respectively, and $r_{ij}$ denotes the known rating of user $i$ on item $j$. We will describe other parameters in details later. 

MF and its variants have been extensively applied to POI recommendation due to their promising performance and scalability \cite{cheng2011exploring,cheng2012fused,yang2013sentiment}. 
However, these methods are all trained by using the centralized mechanism. 
This centralized MF training results in expensive resources required, low model training efficiency, and shallow protection of user privacy. 


\subsection{Decentralized Learning}
Decentralized learning appears to solve the above problems of centralized learning 
\cite{nedic2009distributed,yan2013distributed}. 
Recently, it has been applied in many scenarios such as multiarmed bandit \cite{kalathil2014decentralized}, network distance prediction \cite{liao2010network}, hash function learning \cite{leng2015hashing}, and deep networks \cite{shokri2015privacy,mcmahan2016communication}. 

The most similar existing works to ours are decentralized matrix completion \cite{ling2012decentralized,yun2014nomad}. 
However, we summarize the following two major differences. 
(1) They either allow each learner (user) to communicate with those who have rated the same items or communicate with all the learners, and thus have low accuracy or high communication cost. 
%The neighbors are randomly selected, which can not be applied to real recommendation scenarios. 
In practice, users always collaborate with their affinitive users. 
To capture this, in this paper, we propose a random walk approach for users from the adjacent graph to collaboratively commununicate with each other. 
%we first present an observation by anslysing the data in POI scenarios, i.e., location aggregation. 
%We then propose to build user-user network based on their geographic distance, i.e, the closer they are in geography, the closer they are in the network. 
(2) They allow directly exchange of item preference among learners, which may cause information leakage. 
For example, it is easy to be hacked by using the idea of collaborative filtering \cite{sarwar2001item}, i.e., similar items tend to be preferred by similar users. 
%Assume user $i$ is a malicious user, and he gets items' preferences from user $j'$. If $i$ likes item $j$, and the preference of $j$ and $j'$ are similar, then $i$ will know $i'$ likes $j'$ as well. 
Assume user $i$ is a malicious user, he has his own latent factor of item $j$ ($v^i_j$). 
He also gets the latent factor of item $j$ from user $i'$ ($v^{i'}_j$). 
If $i$ likes item $j$, and $v^i_j$ and $v^{i'}_j$ are similar, then $i$ will know $i'$ likes $j$ as well.
In contrast, in this paper, we propose a gradient exchange scheme to limit the possibility of privacy leakage. 


